% !TEX root = ../main.tex
%% ============================================================================
\section{Related Work and Gaps}
\label{sec:related}
%% ============================================================================

%Two lines of work meet in this paper: assurance of Big Data systems, and automated testing of REST APIs.

%\paragraph{Assurance of Big Data systems.}
The trustworthiness of Big Data computations has been addressed as a property of the pipeline, of the ecosystem of services executing it, and of the data it governs.
Anisetti et al. propose an assurance process that jointly evaluates a Big Data pipeline and its underlying ecosystem, annotating pipeline templates with non-functional requirements and computing an aggregate trustworthiness confidence level, validated on a Big-Data-Analytics-as-a-Service platform~\cite{anisettiAssuranceProcessBig2023}.
Ardagna et al. address the same problem through service-level agreements, negotiating and monitoring contractual guarantees, such as data integrity and availability, between the parties of a Big-Data-Analytics-as-a-Service deployment~\cite{ardagnaBigDataAssurance2023}.
Continuous verification across the pipeline life cycle has been framed as a DevSecOps activity, embedding assurance checks for non-functional requirements at every stage from planning to operation~\cite{anisettiDevSecOpsbasedAssuranceProcess2022}.
At the data governance level, access control has been enforced at ingestion time through data annotations and secure transformations~\cite{anisettiDynamicScalableEnforcement2021}, an approach recently extended into an attribute-based data engine that balances data protection against analytics quality~\cite{polimenoBalancingProtectionQuality2025}.
These contributions target the pipeline, the ecosystem executing it, and the data governed within it; none addresses the REST interface through which every access to that pipeline and that data actually takes place, an assumption on which their own guarantees implicitly rely.
That assumption is what we make verifiable.

%\paragraph{Automated testing of REST APIs.}
%A systematic review of 92 contributions published between 2009 and 2022~\cite{golmohammadiTestingRESTfulAPIs2023} shows that the field concentrated on system testing, understood as verifying that the API responds correctly and respects the declared formats: 72 of the 92 papers fall in this category.
%Security testing is addressed explicitly by 8 papers and listed as an open challenge by 10 more.
Considering the automated testing of REST APIs, it still in its infancy and is considered as an open challenge. Du et al. identified that just 2--5\% of the detected HTTP~500 faults correspond to actual security vulnerabilities~\cite{duVulnerabilityorientedTestingRESTful2024}.

The contract-driven line of research is focused on replacing syntactic heuristics with knowledge of the formal specification.
RESTler infers producer, consumer dependencies from the OpenAPI specification, and generates stateful request sequences. Its oracle is the HTTP~500 status code, which the authors acknowledge cannot reach vulnerabilities that do not manifest through status codes~\cite{atlidakisRESTlerStatefulREST2019a}.
Atlidakis et al. extend it with four formal security rules and active property checkers~\cite{atlidakisCheckingSecurityProperties2020}; NAUTILUS orients the paradigm towards injection-class vulnerabilities requiring ordered operation sequences~\cite{dengNAUTILUSAutomatedRESTful2023}; Schemathesis derives a structured set of conformance properties from the specification and from RFC~7231~\cite{schemathesisDerivingSemanticsawareFuzzers2022}.
A distinct direction automates the construction of the oracle itself, dynamically from observed request/response pairs~\cite{alonsoAGORAAutomatedGeneration2023, alonsoAgora+TestOracleGeneration2025} or statically from the specification~\cite{alonsoSATORIStaticTest2025}; both operate in the domain of functional conformance.
Closest to our objective, Sahin et al. propose nine automated oracles for authorization, authentication, and injection faults, applied as a post-processing phase over an existing fuzzer~\cite{sahinEnhancingRESTAPI2026}; the approach is confined to the application layer and does not cover the configuration of the gateway that fronts the service.
Table~\ref{tab:related} summarises how these approaches use the specification and where their oracle stops.

\begin{table}
  \small
  \caption{What each contract-driven approach checks, and where its oracle stops.}
  \label{tab:related}
  \begin{tabular}{l p{0.28\textwidth} p{0.30\textwidth}}
    \toprule
    \textbf{Approach}                                                                           & \textbf{What it checks}                                                             & \textbf{Oracle boundary}              \\
    \midrule
    RESTler~\cite{atlidakisRESTlerStatefulREST2019a}                                            & Server crashes                                                                      & Only HTTP~500; no semantic check      \\
    Atlidakis et al.~\cite{atlidakisCheckingSecurityProperties2020}                             & Use-after-free, resource leak, hierarchy and namespace violations                   & Fixed rules, no config.\ adaptability \\
    NAUTILUS~\cite{dengNAUTILUSAutomatedRESTful2023}                                            & SQL/command injection, XSS, privilege escalation                                    & Injection-focused only                \\
    Schemathesis~\cite{schemathesisDerivingSemanticsawareFuzzers2022}                           & Schema conformance, undeclared status codes                                         & Structural, no security check         \\
    AGORA/AGORA+~\cite{alonsoAGORAAutomatedGeneration2023,alonsoAgora+TestOracleGeneration2025} & Functional conformance (mined invariants)                                           & Contract only, no security            \\
    SATORI~\cite{alonsoSATORIStaticTest2025}                                                    & Functional conformance (static, via LLM)                                            & Contract only, no security            \\
    Sahin et al.~\cite{sahinEnhancingRESTAPI2026}                                               & Authorization/authentication violations, exposed stack traces, undeclared endpoints & Application layer only, no gateway    \\
    \bottomrule
  \end{tabular}
\end{table}

%The approaches above already expose two of the four gaps that motivate this work.
Given the current related work the following gaps emerge:
Syntactic blindness (\textbf{G1}): what each tool checks, summarised in Table~\ref{tab:related}, is either a status code, a structural property of the response, or a fault class mined or authored independently of the target's security semantics; none of them evaluates whether a syntactically valid request violates the access rules the system is supposed to enforce.
The absence of evaluation criteria for security guarantees (\textbf{G2}): deciding automatically whether an observed behaviour violates such a guarantee requires a criterion that the interface specification alone cannot supply, since it depends on the domain knowledge of the control being exercised.
Barr et al. survey this as the oracle problem, and identify the automated construction of such criteria as the bottleneck of test automation~\cite{barrOracleProblemSoftware2015}. AGORA/AGORA+ and SATORI build oracles explicitly, but mine or infer them from observed traffic or from the contract alone, targeting functional conformance rather than security semantics. Sahin et al. come closer, authoring nine fixed oracles for authorization, authentication, and injection faults, but the result is a closed set confined to the application layer, not a methodology that extends to new domains or to the gateway's own configuration \cite{sahinEnhancingRESTAPI2026}.

\revMin{Lack of Security Contract Validation in Cross-Target Tools (\textbf{G3}): Portable tools can leverage specifications to generate syntactically valid requests across different targets, but they do not verify whether the corresponding responses comply with the security properties declared in those specifications. As a result, portability is achieved at the expense of security-aware semantic validation.}
\revFix{Specification-driven approaches achieve their structural coverage by operating on the syntax of the OpenAPI document, and this same reliance leaves them unable to capture dynamic authorization rules or to detect violations that require understanding the security policies governing resource access}~\cite{abinayaSecurityTestingRESTful2027}.
% Dal paper abinayaSecurityTestingRESTful2027: 
% Although specification-based approaches provide strong structural coverage and automated input generation, they remain constrained by incomplete specifications and limited semantic understanding. Consequently, they often fail to detect undocumented runtime behaviors, authorization flaws, and complex business logic vulnerabilities.
% Dal papaer: Challenges: Specification-driven fuzzing techniques depend on static OAS documents and focus on the syntactic behavior of APIs [37, 57,58]. These methods cannot capture dynamic authorization rules, like role-based restrictions, token reuse, or session expiration. Functional correctness can be validated through status codes and schema vali dation [58,71]. However, detecting authorization violations requires understanding security policies that govern resource access [94].

\revMin{Absence of Deterministic Execution for Continuous Verification (\textbf{G4}): Existing approaches lack deterministic execution, leading to variations in results across successive runs. This prevents reliable comparison of outcomes over time and hinders the integration of automated verification into continuous release and deployment pipelines.}
\revFix{Non-deterministic, or flaky, test outcomes undermine continuous integration by causing builds to fail independently of any genuine fault, an effect observed for automatically generated test suites as much as for developer-written ones}~\cite{parrySurveyFlakyTests2021}. \revFix{An industrial account of API test automation for continuous delivery reports the same symptom directly, noting that automated test suites tend to pass and fail unpredictably, so that even a successful run does not provide reliable confidence to stakeholders}~\cite{bennettPracticalMethodAPI2021}.
% Dal paper parrySurveyFlakyTests2021:
% Tests that fail inconsistently, without changes to the code under test, are described as flaky. Flaky tests do not give a clear indication of the presence of software bugs and thus limit the reliability of the test suites that contain them. A recent survey of software developers found that 59% claimed to deal with flaky tests on a monthly, weekly, or daily basis.
% In a continuous integration system, when a developer has been working on their own branch and goes to merge their changes, their code is remotely built and a test suite run is triggered to ensure their changes have not introduced bugs. Should any tests fail, their changes will be rejected. Naturally, given the association between flaky tests and false alarm failures , flaky tests can limit the efficiency of the continuous integration process by incorrectly causing builds to fail
% [Sez. 4.3.3] These findings demonstrate that automated tools, as well as developers, are also capable of producing flaky tests, threatening the overall reliability of automatically generated test suites.
% [Sez. 4.3.3] regression test suites are useful in as far as they capture the behavior of a program at a given point in time and are used to identify if a recent code change has any unintended effects. If a regression test suite contains flaky tests, then it does not accurately reflect the behavior of the program and is thus less informative for developers.

% Dal paper bennettPracticalMethodAPI2021:
% The first problem is that execution of tests tends to fail and pass in seemingly nondeterministic ways (tests are flaky).
% The automated tests failed and passed in seemingly nondeterministic ways - the tests were flaky.
% So even a successful test run did not provide clear confidence for relevant stakeholders.


